% TOP_PROBLEM: {"problemId": "problem.qma-versus-qcma-problem", "canonicalTitle": "QMA versus QCMA Problem", "releaseRank": 122, "primaryDomain": "quantum_information_computation", "primaryDomainLabel": "Quantum information & computation", "displayStatus": "open_with_solved_subcases", "releaseStatus": "open_with_solved_subcases"}
% !TeX program = lualatex
% Definition and sources only; this file is not an LLM solution attempt.
\documentclass[11pt]{article}
\usepackage[margin=25mm]{geometry}
\usepackage{fontspec}
\setmainfont{Latin Modern Roman}
\usepackage{amsmath,amssymb,mathtools}
\usepackage{unicode-math}
\setmathfont{Latin Modern Math}
\usepackage{xurl,hyperref}
\hypersetup{colorlinks=true,urlcolor=blue}
\setcounter{secnumdepth}{0}
\setlength{\parindent}{0pt}
\setlength{\parskip}{5pt}
\setlength{\emergencystretch}{3em}
\begin{document}
\section*{\texorpdfstring{QMA versus QCMA Problem}{QMA versus QCMA Problem}}\label{122}
\subsection{Definitions and mathematical statement}
A QMA verifier is a uniform polynomial-time quantum computation receiving a polynomial-size quantum witness. A QCMA verifier has the same quantum computational power but receives only a polynomial-length classical bit string. Both use completeness at least $2/3$ and soundness at most $1/3$, with soundness quantified over every allowed witness on no-inputs. The question is
\[
 \mathsf{QMA}\stackrel{?}{=}\mathsf{QCMA}.
\]
The catalogue states a language formulation; a promise-problem version should be labeled separately rather than silently substituted. An oracle separation does not determine the ordinary unrelativized relation.
\par\smallskip\textit{Formulation and concept references:} \href{https://eccc.weizmann.ac.il/report/2025/176/revision/1/download/}{[S1]}.

\subsection{Short English statement}
Can quantum evidence always be replaced by classical evidence without sacrificing efficient quantum verification?
\subsection{Sources}
\begin{itemize}
\item[S1]John Bostanci, Jonas Haferkamp, Chinmay Nirkhe, and Mark Zhandry, Separating QMA from QCMA with a classical oracle, Introduction.
\url{https://eccc.weizmann.ac.il/report/2025/176/revision/1/download/}.
\textit{Formulation source.}
\end{itemize}
\end{document}
